\chapter{Adelic Categorical Descent}
\label{v4-arith:adelic-descent}

\(\mathrm{L}_{\mathrm{ACD}}\) (\Cref{v4-arith:cl:cor:F1-single-lemma})
is the essential adelic gluing statement for conjecture F1 (global
arithmetic Arinkin--Gaitsgory for \(GL_{2}/\mathbb{Q}\)), with
a substantial closure problem
(App.~\ref{v4-arith:app:three-conjectures-dep}). This chapter decomposes
\(\mathrm{L}_{\mathrm{ACD}}\) into four subconditions
\((L_{1})\)--\((L_{4})\).  The formal restricted tensor product
\((L_{1})\) is constructed here.  The archimedean representation
category \((L_{4})\) is a boundary datum, not a spectral equivalence.
\((L_{2})\) still needs the local categorical comparison and the
Whittaker compatibility.  \((L_{3})\) is the Bernstein-centre
diamond-gluing compatibility on local coherent Springer sheaves over
the Selmer-derived global spectral stack.  The main result
(\Cref{v4-arith:acd:thm:final-reduction}) is a conditional reduction,
not a proof of \(\mathrm{L}_{\mathrm{ACD}}\).

\section{Statement Recalled}
\label{v4-arith:acd:statement}

\subsection{The Condition}

\Cref{v4-arith:cl:cor:F1-single-lemma} states:

\begin{quote}
\textbf{\(\mathrm{L}_{\mathrm{ACD}}\) (adelic categorical descent).}
The assignment
\begin{equation}
\label{v4-arith:acd:eq:assignment}
v\;\longmapsto\;\mathrm{D}^{\mathrm{sph\text{-}fin}}\bigl([GL_{2}(\mathbb{Q}_{v})]\bigr)
\end{equation}
(for \(v\) ranging over places of \(\mathbb{Q}\), finite and archimedean)
forms a Tate-restricted tensor product whose colimit is the global
automorphic category
\(\mathrm{D}([GL_{2}])^{\mathrm{sph\text{-}fin}}\); this restricted
tensor product is compatible with Hecke-functoriality and with the
local Fargues--Scholze spectral action together with the still-open
F1.a local categorical comparison, so that
\begin{equation}
\label{v4-arith:acd:eq:target}
\bigotimes_{v}^{\prime}\mathrm{D}^{\mathrm{sph\text{-}fin}}\bigl([GL_{2}(\mathbb{Q}_{v})]\bigr)
\;\simeq\;
\mathrm{D}\bigl([GL_{2}]\bigr)^{\mathrm{sph\text{-}fin}}
\end{equation}
as Hecke-equivariant stable presentable \(\infty\)-categories.
\end{quote}

The restricted tensor product \(\bigotimes'_{v}\) follows Tate 1950 /
Weil 1967 at the level of spaces of automorphic forms. The canonical
\(\infty\)-categorical definition is the filtered-colimit construction
of \Cref{v4-arith:tate:def:restricted-tensor}. The display below is
mnemonic:
\begin{equation}
\label{v4-arith:acd:eq:restricted-tp}
\bigotimes_{v}^{\prime}\mathcal{C}_{v}
\;:=\;
\varinjlim_{S\text{ finite}}
\bigotimes_{v\in S}\mathcal{C}_{v}\otimes\bigotimes_{v\notin S}\mathbb{1}_{u_{v}},
\end{equation}
where \(u_{v}\in\mathcal{C}_{v}\) is the distinguished unit
(the unit of the unramified spherical Hecke category) at almost every
finite place.

\subsection{Context and Significance}

\(\mathrm{L}_{\mathrm{ACD}}\) is the essential adelic gluing
statement for F1 (global arithmetic Arinkin--Gaitsgory,
\Cref{v4-arith:conj:F1-global-AG}). The structural reduction of F1
in \Cref{v4-arith:thm:F1-reduction} splits F1 into four sub-items
 (F1.a)--(F1.d). Only (F1.c) is a theorem drawn directly from
 primary literature (Emerton--Gee arXiv:1908.07185). Fargues--Scholze
 arXiv:2102.13459 supplies local geometrisation and spectral action,
 not the categorical equivalence required by (F1.a). The remaining
 items (F1.a), (F1.b), and (F1.d) are equivalent to the local
comparison plus \(\mathrm{L}_{\mathrm{ACD}}\) conditions whose
closure removes F1.

\(\mathrm{L}_{\mathrm{ACD}}\) is a
\textbf{conjectural, non-formal theorem}. This chapter does not
close \(\mathrm{L}_{\mathrm{ACD}}\); it reduces it to a narrower and
precisely-stated compatibility on local coherent Springer sheaves
(Ben-Zvi--Chen--Helm--Nadler arXiv:2010.02321) glued across primes
on the Selmer-derived global spectral stack.

\section{Modular Reduction of \(\mathrm{L}_{\mathrm{ACD}}\)}
\label{v4-arith:acd:obstruction-analysis}

Six independent lines of obstruction on \(\mathrm{L}_{\mathrm{ACD}}\) are
given below. Each line produces a structural decomposition, a partial
closure, or a named residual, following the schema of
\Cref{v4-arith:cl:criterion}.

\subsection{Modular Decomposition of \(\mathrm{L}_{\mathrm{ACD}}\)}
\label{v4-arith:acd:obstruction-1}

\(\mathrm{L}_{\mathrm{ACD}}\) is the conjunction of four ingredients:
(i) well-definedness of the restricted tensor product of local
\(\infty\)-categories with distinguished units;
(ii) Hecke-functoriality across places;
(iii) compatibility of the Bernstein centre action across places
under restricted tensor;
(iv) archimedean-finite assembly (extension across the archimedean
place).
A structural closure requires separating each ingredient. The
decomposition admits four disjoint subconditions,
\((L_{1})\)--\((L_{4})\), each a precise statement at the
\(\infty\)-categorical level.

\begin{theorem}[modular decomposition of \(\mathrm{L}_{\mathrm{ACD}}\)]
\label{v4-arith:acd:thm:modular-decomposition}
\ClaimStatusProvedHere. \(\mathrm{L}_{\mathrm{ACD}}\) contains the
following four subconditions:
\begin{itemize}
\item[\((L_{1})\)] \emph{Well-definedness of the restricted tensor
product.} The assignment \(v\mapsto\mathcal{C}_{v}\) with
\(\mathcal{C}_{v}=\mathrm{D}^{\mathrm{sph\text{-}fin}}([GL_{2}(\mathbb{Q}_{v})])\)
and distinguished units \(u_{v}\) (spherical vectors) defines a
well-formed restricted tensor product \(\bigotimes'_{v}\mathcal{C}_{v}\)
in the \(\infty\)-category \(\mathrm{Pr}^{\mathrm{L,st}}\) of stable
presentable \(\infty\)-categories.
\item[\((L_{2})\)] \emph{Hecke-functoriality lift across places.} For
each finite set \(S\) of places, assume the local categorical
comparison functors
\(\mathcal{C}_{v}\to\mathrm{IndCoh}_{\mathrm{Nilp}}(\mathcal{X}^{\mathrm{EG},v}_{GL_{2}})\)
for the finite places \(v\in S\).  Their tensor product is required to
give an equivalence
\(\bigotimes_{v\in S}\mathcal{C}_{v}\simeq\mathrm{IndCoh}_{\mathrm{Nilp}}(\prod_{v\in S}\mathcal{X}^{\mathrm{EG},v}_{GL_{2}})\)
compatible with Hecke operators at distinct primes.
\item[\((L_{3})\)] \emph{Bernstein-centre diamond gluing.} The
Bernstein-centre action
\(\mathcal{Z}_{v}\curvearrowright\mathcal{C}_{v}\) (the centre of the
local categorical Hecke algebra) extends to a global action
\(\mathcal{Z}_{\mathrm{glob}}\curvearrowright\bigotimes'_{v}\mathcal{C}_{v}\)
making the tensor product restricted-tensor-equivariant under the
Bernstein centre.
\item[\((L_{4})\)] \emph{Archimedean extension.} The archimedean
factor \(\mathcal{C}_{\infty}=\mathrm{D}^{\mathrm{sph\text{-}fin}}([GL_{2}(\mathbb{R})])\)
is the ind-completed derived category of admissible
\((\mathfrak{gl}_{2},O(2))\)-modules of finite \(K\)-type.  It carries
the trivial \(K\)-spherical object as distinguished unit, and the
restricted tensor product
\(\bigotimes'_{v}\mathcal{C}_{v}\) extends the finite-place product
\(\bigotimes'_{p<\infty}\mathcal{C}_{p}\) by tensoring with
\(\mathcal{C}_{\infty}\).
\end{itemize}
\(\mathrm{L}_{\mathrm{ACD}}\) holds if \((L_{1})\) through
\((L_{4})\) hold together with the identification of the resulting
object with the global automorphic category
\(\mathrm{D}([GL_{2}])^{\mathrm{sph\text{-}fin}}\).
\end{theorem}

\begin{proof}
Forward direction: assume \(\mathrm{L}_{\mathrm{ACD}}\). The restricted
tensor product~\eqref{v4-arith:acd:eq:restricted-tp} is well-defined
by definition, giving \((L_{1})\). Hecke-functoriality across places
is part of the \(\mathrm{L}_{\mathrm{ACD}}\) hypothesis, giving
\((L_{2})\). Bernstein-centre compatibility follows from functoriality
of the global automorphic category under the global Bernstein centre
plus Hecke-functoriality, giving \((L_{3})\). Archimedean extension is
explicit in the definition of \(\mathrm{L}_{\mathrm{ACD}}\) (the index
\(v\) ranges over all places), giving \((L_{4})\).

Converse direction: assume \((L_{1})\)--\((L_{4})\) and the global
identification just stated. By \((L_{1})\) and
\((L_{4})\), \(\bigotimes'_{v}\mathcal{C}_{v}\) is a well-defined stable
presentable \(\infty\)-category. By \((L_{2})\) and \((L_{3})\), this
category is Hecke-equivariant and Bernstein-centre-equivariant. The
identification with the global automorphic category follows from
the assumed local-global compatibility. No theorem of Bernstein,
Raskin, or geometric Langlands supplies this number-field
identification without that input.
\end{proof}

\begin{remark}[form of subconditions at statement]
\label{v4-arith:acd:rem:subcondition-form}
At statement (2026-04-24) the subconditions carry the following
form:
\begin{description}
\item[\((L_{1})\)] Theorem, closed in this chapter
(\S\ref{v4-arith:acd:L1}).
\item[\((L_{2})\)] Conditional; reduces to the local categorical
comparison plus Whittaker and Bernstein-centre compatibility
(\S\ref{v4-arith:acd:L2}).
\item[\((L_{3})\)] Conjectural; the residual content of
\(\mathrm{L}_{\mathrm{ACD}}\) after the decomposition of
\Cref{v4-arith:acd:thm:modular-decomposition}
(\S\ref{v4-arith:acd:L3}).
\item[\((L_{4})\)] Boundary datum; Jacquet--Langlands and Vogan define
the archimedean representation category, not an Emerton--Gee spectral
equivalence (\S\ref{v4-arith:acd:L4}).
\end{description}
\end{remark}

\subsection{Curve Precedent and Transport Obstructions}
\label{v4-arith:acd:obstruction-2}

For a smooth complete curve over an algebraically closed field of
characteristic zero, the de Rham geometric Langlands theorem gives the
model for the categorical descent pattern.  The sheaf-function
analogue over finite fields is a different theorem.  Neither theorem
transports directly to the number-field case.  The obstructions are:
\begin{enumerate}
\item No archimedean place in the curve setting. Geometric Langlands
uses the closed points of a curve. Number-field \(\mathrm{L}_{\mathrm{ACD}}\)
requires the archimedean extension \((L_{4})\), which has no curve
analogue.
\item Emerton--Gee integrality. The number-field spectral side uses
the Emerton--Gee stack arXiv:1908.07185, which encodes integrality
of Galois representations; de Rham geometric Langlands uses the stack
\(\mathrm{LocSys}^{\mathrm{dR}}_{G}(X)\), which is smoother and has
weaker integrality content. The Bernstein-centre compatibility
\((L_{3})\) is strictly more subtle in the number-field setting
because the Emerton--Gee stack is not smooth over
\(\mathrm{Spec}\,\mathbb{Z}\).
\item Perfectoid diamonds. Fargues--Scholze's local theorem uses the
Fargues--Fontaine curve and perfectoid diamonds; geometric Langlands
uses ind-schemes. Gluing across finite places in the
number-field case requires comparing perfectoid diamonds at
distinct primes, a non-issue in the de Rham curve setting.
\end{enumerate}

\begin{theorem}[curve transport, qualified]
\label{v4-arith:acd:thm:function-field-transport}
\ClaimStatusNeedsVerification. The Gaitsgory--Rozenblyum--Raskin--Lin /
AGKRRV proof of de Rham geometric Langlands gives the model for
subconditions \((L_{1})\) and \((L_{2})\) as an analogy, not as a
transport theorem (Raskin,
``Chiral principal series categories I'', Adv. Math. 388 (2021);
Raskin, ``D-modules on infinite-dimensional varieties'', available
as \texttt{samraskin.net/dmod.pdf}). It closes \((L_{1})\) only after
the number-field local categories and distinguished units are fixed.
It does not close \((L_{2})\) without the local categorical comparison
of (F1.a) and (BC-diamond-glue).
It does \emph{not} close \((L_{3})\) (Bernstein centre compatibility
is strictly more subtle in the number-field setting due to
Emerton--Gee stack non-smoothness) or \((L_{4})\) (no archimedean
analogue on the curve).
\end{theorem}

\begin{proof}
\((L_{1})\): Raskin 2021 JAMS constructs restricted tensor products of
local principal-series categories over a curve. The construction is
\(\infty\)-categorical and uses only (a) stability, (b) presentability,
(c) existence of distinguished units. These properties can be imposed
on the number-field local categories. This gives the formal construction
of \((L_{1})\), not the global automorphic identification.

\((L_{2})\): AGKRRV's Hecke-functoriality argument combines local
geometric Langlands equivalences with Bernstein--Zelevinsky
factorization on a curve. For number fields,
Fargues--Scholze arXiv:2102.13459 supplies the local diamond
geometrisation and spectral action. It does not supply the displayed
categorical comparison with the Emerton--Gee stack. The
Bernstein--Zelevinsky pattern applies only after the local
comparison and Whittaker-model compatibility are constructed. Domain:
\S\ref{v4-arith:acd:L2}.

\((L_{3})\): Number-field obstruction. The Bernstein centre
\(\mathcal{Z}^{\mathrm{Ar}}\) of the global arithmetic Iwahori Hecke
category acts on \(\bigotimes'_{v}\mathcal{C}_{v}\) via the local
Bernstein centres \(\mathcal{Z}_{v}\), but compatibility across primes
involves comparing Bernstein parameters at distinct primes through
the Emerton--Gee stack. The curve analogue uses the de~Rham
stack, which is smoother. The number-field compatibility is the
residual.

\((L_{4})\): No curve analogue.
\end{proof}

\subsection{Sphericity-Finiteness Preserved Under Restricted Tensor}
\label{v4-arith:acd:obstruction-3}

The ``sph-fin'' cut in
\(\mathrm{D}^{\mathrm{sph\text{-}fin}}([GL_{2}(\mathbb{Q}_{v})])\)
restricts to categories with finite-type spherical sub-structure,
excluding the continuous Eisenstein spectrum and broad classes of
non-tempered representations. The condition is preserved under the
restricted tensor product.

\begin{lemma}[sph-fin preservation under restricted tensor]
\label{v4-arith:acd:lem:sph-fin-preservation}
\ClaimStatusProvedHere. Let \(\{\mathcal{C}_{v}\}\) be a family of
stable presentable \(\infty\)-categories with distinguished units
\(u_{v}\) at almost every \(v\), and assume each \(\mathcal{C}_{v}\) is
``sph-fin'' in the sense that: (i) \(\mathcal{C}_{v}\) has a compact
subcategory \(\mathcal{C}_{v}^{\omega}\) whose compact objects are
finite-type over the local Hecke algebra \(\mathcal{H}_{v}\); (ii)
\(u_{v}\) is compact in \(\mathcal{C}_{v}^{\omega}\); (iii) the compact
objects generate \(\mathcal{C}_{v}\). Then the restricted tensor
product \(\bigotimes'_{v}\mathcal{C}_{v}\) is likewise sph-fin.
\end{lemma}

\begin{proof}
By Lurie HA 5.3 (higher Deligne centralizers) and HA 4.8 (tensor
products of presentable \(\infty\)-categories), the tensor product
\(\bigotimes_{v\in S}\mathcal{C}_{v}\) over a finite set \(S\) inherits
compact generation from its finite factors: a family of compact
generators is given by tensor products of local compact generators.
The distinguished units \(u_{v}\) at \(v\notin S\) are compact, so
tensoring with them preserves compactness. Taking the filtered
colimit over \(S\) preserves compact generation (Lurie HA 5.3.2.9 on
filtered colimits of compactly generated categories). Hence
\(\bigotimes'_{v}\mathcal{C}_{v}\) is compactly generated, with
compact subcategory the filtered colimit of finite tensor products
of local compacts. The sph-fin condition is preserved.
\end{proof}

\subsection{Cross-Place Hecke Commutativity}
\label{v4-arith:acd:obstruction-4}

Commutativity of \(T_{p}\) and \(T_{q}\) at the level of cohomology is
automatic; categorical commutativity in the \(\infty\)-setting is a
coherence datum following from the \(E_{2}\)-monoidal structure on
each local Hecke category (assuming F2). Cross-place compatibility
requires showing that the global automorphic category decomposes as a
restricted tensor product whose factor at each prime is the local Hecke
category.

\begin{proposition}[cross-place Hecke commutativity]
\label{v4-arith:acd:prop:cross-place-commutativity}
\ClaimStatusProvedHere. In the restricted tensor product
\(\bigotimes'_{v}\mathcal{C}_{v}\), the local Hecke operators
\(T_{p}:\mathcal{C}_{p}\to\mathcal{C}_{p}\) extended by identity on
factors at \(v\neq p\) (Tate extension) commute: for distinct primes
\(p\neq q\),
\begin{equation}
(T_{p}\otimes\mathrm{id}_{\neq p})\circ(T_{q}\otimes\mathrm{id}_{\neq q})\;=\;(T_{q}\otimes\mathrm{id}_{\neq q})\circ(T_{p}\otimes\mathrm{id}_{\neq p})
\end{equation}
as endofunctors of \(\bigotimes'_{v}\mathcal{C}_{v}\), up to canonical
homotopy.
\end{proposition}

\begin{proof}
Each \(T_{p}\) acts on the \(\mathcal{C}_{p}\)-factor and by identity on
all other factors. The universal property of the tensor product in
\(\mathrm{Pr}^{\mathrm{L,st}}\) (Lurie HA 4.8.1) gives that
endofunctors factoring through disjoint tensor factors commute up
to canonical homotopy. Since \(T_{p}\) factors through \(\mathcal{C}_{p}\)
and \(T_{q}\) through \(\mathcal{C}_{q}\), their extensions to the
restricted tensor product commute.
\end{proof}

\subsection{Non-Transference from Caraiani--Scholze}
\label{v4-arith:acd:obstruction-5}

Caraiani--Scholze arXiv:1511.02418 prove a unitary Shimura variety
analogue of \(\mathrm{L}_{\mathrm{ACD}}\) for signature \((1,n-1)\)
unitary groups \(G=U(1,n-1)\) over CM fields. A Jacquet--Langlands
transfer from \(GL_{2}/\mathbb{Q}\) to \(U(1,1)\) over an imaginary
quadratic field \(F/\mathbb{Q}\) would reduce the problem to the
Caraiani--Scholze case; this reduction fails. Caraiani--Scholze's proof
uses three inputs specific to unitary Shimura varieties:
\begin{enumerate}
\item Calegari--Geraghty patching arXiv:1207.4224 for the
Galois-deformation side, producing a patched space of automorphic
forms whose properties reflect the Emerton--Gee stack.
\item The Kottwitz conjecture on cohomology of Shimura varieties in
its form for unitary groups (Caraiani--Scholze 2017 \S4).
\item Ihara-avoidance hypothesis on the existence of adequate
Taylor--Wiles primes, required for patching.
\end{enumerate}

None of these transfer to \(GL_{2}/\mathbb{Q}\) in the generality
\(\mathrm{L}_{\mathrm{ACD}}\) requires.

\begin{observation}[non-transference of Caraiani--Scholze]
\label{v4-arith:acd:obs:non-transference-CS}
\(GL_{2}/\mathbb{Q}\) has a Shimura variety over \(\mathbb{Q}\) in the
form of the modular curve \(Y(\Gamma(N))\), but the Shimura variety is
\(1\)-dimensional, whereas Caraiani--Scholze's Ihara-avoidance argument
requires the Shimura variety to be higher-dimensional (specifically
\(\geq n-1\) where \(n=\mathrm{rk}\,G\)). For \(GL_{2}/\mathbb{Q}\), the
Ihara-avoidance input is essentially degenerate. Caraiani--Scholze's
proof cannot be transferred.
\end{observation}

The transference fails; but Caraiani--Scholze's proof does suggest
the structural shape \(\mathrm{L}_{\mathrm{ACD}}\) must take: a
patched-spaces argument with local Fargues--Scholze input glued
across primes by an ``Ihara-avoidance-free'' version of the
argument.

\subsection{Reduction to a Single Compatibility}
\label{v4-arith:acd:obstruction-6}

\((L_{1})\) is formal in presentable categories.  \((L_{4})\) supplies
the archimedean representation-theoretic boundary datum.
\((L_{2})\) reduces to the local categorical comparison plus
Whittaker-model compatibility; the remaining central compatibility is
\((L_{3})\): Bernstein-centre diamond gluing.
\(\mathrm{L}_{\mathrm{ACD}}\) therefore reduces to the single
compatibility stated below.

\begin{theorem}[conditional reduction of \(\mathrm{L}_{\mathrm{ACD}}\)]
\label{v4-arith:acd:thm:final-reduction}
\ClaimStatusProvedHereConditional. Assuming the local categorical
comparison at all finite primes and the archimedean boundary datum,
\(\mathrm{L}_{\mathrm{ACD}}\) reduces to the following compatibility:
\begin{itemize}
\item[\textbf{(BC-diamond-glue)}] \emph{Bernstein-centre diamond
gluing.} For each pair of primes \(p\neq q\), the local Bernstein
centres \(\mathcal{Z}_{p}\) and \(\mathcal{Z}_{q}\) (the centres of the
local categorical Hecke algebras on the perfectoid Fargues--Fontaine
curves at \(p\) and \(q\)) admit a common global refinement
\(\mathcal{Z}_{\{p,q\}}\) such that:
\begin{itemize}
\item[(a)] \(\mathcal{Z}_{\{p,q\}}\) embeds into
\(\mathcal{Z}_{p}\otimes\mathcal{Z}_{q}\) under the external tensor
action on \(\mathcal{C}_{p}\otimes\mathcal{C}_{q}\);
\item[(b)] the embedding is compatible with restriction to smaller
prime sets \(S\subsetneq\{p,q\}\);
\item[(c)] the inductive limit
\(\mathcal{Z}_{\mathrm{glob}}:=\varinjlim_{S}\mathcal{Z}_{S}\) over
finite sets \(S\) of primes agrees with the global Bernstein centre
acting on the global automorphic category.
\end{itemize}
\end{itemize}
\end{theorem}

\begin{proof}
Forward direction: assume \(\mathrm{L}_{\mathrm{ACD}}\). The global
automorphic category \(\mathrm{D}([GL_{2}])^{\mathrm{sph\text{-}fin}}\)
carries a global Bernstein-centre action. The restricted tensor
product structure of \(\mathrm{L}_{\mathrm{ACD}}\) expresses this
global action as a filtered colimit of finite products of local
Bernstein centres. Compatibilities (a)--(c) follow.

Reverse direction: assume (BC-diamond-glue), the local categorical
comparisons, and the archimedean boundary datum. Then by
\Cref{v4-arith:acd:thm:modular-decomposition}, it suffices to close
\((L_{1})\), \((L_{2})\), \((L_{3})\), \((L_{4})\). \((L_{1})\) closes
in~\S\ref{v4-arith:acd:L1}. \((L_{2})\) is supplied
by the assumed local comparisons, up to local-global compatibility of
Whittaker models, which (BC-diamond-glue) supplies via
compatibility~(b). \((L_{3})\) closes from (BC-diamond-glue) directly.
\((L_{4})\) is the assumed archimedean boundary datum. Assembling,
\(\mathrm{L}_{\mathrm{ACD}}\) holds.
\end{proof}

\begin{remark}[domain of (BC-diamond-glue)]
\label{v4-arith:acd:rem:bc-diamond-glue-domain}
(BC-diamond-glue) is strictly narrower than \(\mathrm{L}_{\mathrm{ACD}}\).
It is a single compatibility between local Bernstein centres and the
global Bernstein centre. Partial change:
Ben-Zvi--Chen--Helm--Nadler arXiv:2010.02321 (``Coherent Springer
theory and the categorical Deligne--Langlands correspondence'')
constructs the local Bernstein centres as coherent sheaves on the
local Emerton--Gee stacks. (BC-diamond-glue) asks that these coherent
sheaves glue across primes to a global coherent sheaf on the
Selmer-derived global spectral stack
\(\mathcal{X}^{\mathrm{Sel}}_{\bar\rho,\det,\mathcal{L}}\), agreeing
with the global Bernstein centre action at the level of categorical
trace.
This is conjectural but rests on concrete partial change.

Required theorem of (BC-diamond-glue): non-formal theorem, as opposed
to the non-formal theorem condition for \(\mathrm{L}_{\mathrm{ACD}}\).
The reduction from
\Cref{v4-arith:acd:thm:final-reduction} is a factor of two to three
in domain.
\end{remark}

\section{Subcondition \((L_{1})\): Restricted Tensor in \(\infty\)-Categories}
\label{v4-arith:acd:L1}

\subsection{Definition}

\begin{definition}[restricted tensor product in \(\mathrm{Pr}^{\mathrm{L,st}}\)]
\label{v4-arith:acd:def:restricted-tensor}
The canonical definition is \Cref{v4-arith:tate:def:restricted-tensor};
the following formula is its older mnemonic form in the ACD notation.
Let \(V\) be a countable indexing set, \(\{\mathcal{C}_{v}\}_{v\in V}\)
a family of stable presentable \(\infty\)-categories, and let
\(u_{v}\in\mathcal{C}_{v}\) be a distinguished unit at almost every
\(v\) (i.e.\ at all but finitely many \(v\) in a fixed cofinite subset
\(V^{\mathrm{unr}}\subseteq V\)). The \emph{Tate-restricted tensor
product} is
\begin{equation}
\bigotimes_{v\in V}^{\prime}\mathcal{C}_{v}\;:=\;\varinjlim_{S\text{ finite},\; V^{\mathrm{unr}}\supseteq V\setminus S}\;\bigotimes_{v\in S}\mathcal{C}_{v}\otimes\bigotimes_{v\in V^{\mathrm{unr}}\setminus S}\mathbb{1}_{u_{v}},
\end{equation}
where \(\mathbb{1}_{u_{v}}\) denotes the unit \(\infty\)-category
generated by \(u_{v}\) (a full subcategory of \(\mathcal{C}_{v}\)).
\end{definition}

\begin{theorem}[\((L_{1})\) closed]
\label{v4-arith:acd:thm:L1-closed}
\ClaimStatusProvedHere. The restricted tensor product of
\Cref{v4-arith:acd:def:restricted-tensor} is a well-defined stable
presentable \(\infty\)-category. Specifically:
\begin{enumerate}
\item The filtered colimit \(\varinjlim_{S}\) exists in
\(\mathrm{Pr}^{\mathrm{L,st}}\).
\item The resulting object is stable, presentable, and
compactly-generated if each \(\mathcal{C}_{v}\) is compactly-generated
(Raskin, ``D-modules on infinite-dimensional varieties'',
\texttt{samraskin.net/dmod.pdf}; Lurie HA 5.3.2).
\item Sphericity-finiteness is preserved
(\Cref{v4-arith:acd:lem:sph-fin-preservation}).
\item The restricted tensor product is functorial in the family
\(\{\mathcal{C}_{v}, u_{v}\}\) under morphisms that preserve
distinguished units.
\end{enumerate}
\end{theorem}

\begin{proof}
\textbf{(i) Existence of the filtered colimit.} The \(\infty\)-category
\(\mathrm{Pr}^{\mathrm{L,st}}\) of stable presentable \(\infty\)-categories
is a (large) \(\infty\)-category admitting all small colimits (Lurie
HA 4.8.1.15). Filtered colimits exist.

\textbf{(ii) Presentability and compact generation.} Lurie HA
5.3.2.9 asserts that a filtered colimit of compactly-generated
stable presentable \(\infty\)-categories, under functors preserving
compact objects, is compactly generated. The transition maps in the
defining filtered colimit are of the form ``tensor with
\(\mathbb{1}_{u_{v}}\)'' for newly-included places \(v\), which are
left-adjoint (and preserve compact objects because the unit \(u_{v}\)
is assumed compact, as in
\Cref{v4-arith:acd:lem:sph-fin-preservation}). Hence the colimit is
compactly generated.

\textbf{(iii) Stability.} Filtered colimits of stable
\(\infty\)-categories in \(\mathrm{Pr}^{\mathrm{L,st}}\) are stable
(Lurie HA 1.1.4.4: stable \(\infty\)-categories are closed under
filtered colimits in \(\mathrm{Pr}^{\mathrm{L}}\)).

\textbf{(iv) Sphericity-finiteness preservation.}
\Cref{v4-arith:acd:lem:sph-fin-preservation}.

\textbf{(v) Functoriality.} A morphism of families
\(\{(\mathcal{C}_{v}, u_{v})\}\to\{(\mathcal{C}'_{v}, u'_{v})\}\) is a
family of left adjoints \(F_{v}:\mathcal{C}_{v}\to\mathcal{C}'_{v}\)
sending \(u_{v}\) to \(u'_{v}\). The induced map
\(\bigotimes_{v\in S}F_{v}:\bigotimes_{v\in S}\mathcal{C}_{v}\to\bigotimes_{v\in S}\mathcal{C}'_{v}\)
is compatible with unit-extension to larger \(S'\supseteq S\) because
\(F_{v}\) sends distinguished units to distinguished units. Taking the
filtered colimit yields the induced map
\(\bigotimes'_{v}F_{v}:\bigotimes'_{v}\mathcal{C}_{v}\to\bigotimes'_{v}\mathcal{C}'_{v}\).
The construction is independent of the choice of cofinite unramified
subset \(V^{\mathrm{unr}}\) (any two choices agree on a further
cofinite subset, and the filtered colimit is insensitive to such
shrinking).
\end{proof}

\begin{remark}[comparison with Raskin's construction]
\label{v4-arith:acd:rem:raskin-comparison}
Raskin ``Chiral principal series categories I'', Adv. Math. 388
(2021), 107856, and Raskin ``D-modules on infinite-dimensional
varieties'' (\texttt{samraskin.net/dmod.pdf}) construct
restricted tensor products of principal-series categories over a
smooth curve. The construction uses the same filtered colimit
pattern with spherical vectors as distinguished units. The
number-field analogue is not a direct transport.  The same filtered
colimit defines the formal restricted tensor product once the local
categories and distinguished units are fixed.  The non-trivial content
is the identification with the global automorphic category, which is
\((L_{2})\) and \((L_{3})\).
\end{remark}

\begin{corollary}[\((L_{1})\) holds]
\label{v4-arith:acd:cor:L1-holds}
\ClaimStatusProvedHere. \((L_{1})\) is a theorem, closed by
\Cref{v4-arith:acd:thm:L1-closed} plus
\Cref{v4-arith:acd:lem:sph-fin-preservation}.
\end{corollary}

\section{Subcondition \((L_{2})\): Hecke-Functoriality Across Places}
\label{v4-arith:acd:L2}

\subsection{Statement and Reduction}

\begin{theorem}[\((L_{2})\) reduction]
\label{v4-arith:acd:thm:L2-reduction}
\ClaimStatusNeedsVerification. For each finite set \(S\) of primes, the
Fargues--Scholze local geometry, after the local categorical comparison
(F1.a), supplies comparison functors
\begin{equation}
\label{v4-arith:acd:eq:local-FS}
\Phi_{v}^{\mathrm{loc}}\;:\;\mathcal{C}_{v}\;\longrightarrow\;\mathrm{IndCoh}_{\mathrm{Nilp}}\bigl(\mathcal{X}^{\mathrm{EG},v}_{GL_{2}}\bigr),\qquad v\in S
\end{equation}
which glue to an equivalence
\begin{equation}
\label{v4-arith:acd:eq:glued-FS}
\Phi_{S}^{\mathrm{loc}}\;:\;\bigotimes_{v\in S}\mathcal{C}_{v}\;\xrightarrow{\sim}\;\mathrm{IndCoh}_{\mathrm{Nilp}}\biggl(\prod_{v\in S}\mathcal{X}^{\mathrm{EG},v}_{GL_{2}}\biggr)
\end{equation}
compatible with Hecke operators \(T_{v}\) at distinct primes. The
open content is the local categorical comparison plus the
(BC-diamond-glue) Whittaker/Bernstein-centre compatibility.
\end{theorem}

\begin{proof}
\textbf{Step 1: Local geometry exists.} At each prime \(p\),
Fargues--Scholze arXiv:2102.13459 constructs local geometrization,
spectral action, and Fargues--Fontaine geometry. It does not by
itself identify \(\mathcal{C}_{p}\) with
\(\mathrm{IndCoh}_{\mathrm{Nilp}}(\mathcal{X}^{\mathrm{EG},p}_{GL_{2}})\).
That identification is the local comparison input (F1.a). At the
archimedean place, the analogue is Jacquet--Langlands SLN 114
(see~\S\ref{v4-arith:acd:L4}).

\textbf{Step 2: Hecke-functoriality at each place.} Local
Fargues--Scholze includes Hecke-functoriality at each prime: the
local Hecke algebra \(\mathcal{H}_{p}\) acts on both sides of
\(\Phi_{p}^{\mathrm{loc}}\) once the local comparison is supplied, and
\(\Phi_{p}^{\mathrm{loc}}\) must be
\(\mathcal{H}_{p}\)-equivariant.

\textbf{Step 3: Tensor product of equivalences.} Lurie HA 4.8.1
states that the tensor product in \(\mathrm{Pr}^{\mathrm{L,st}}\) is
functorial: a family of equivalences
\(\Phi_{v}:\mathcal{C}_{v}\to\mathcal{D}_{v}\) produces an equivalence
\(\bigotimes_{v\in S}\Phi_{v}:\bigotimes_{v\in S}\mathcal{C}_{v}\to\bigotimes_{v\in S}\mathcal{D}_{v}\).
Applied to \(\Phi_{v}^{\mathrm{loc}}\) (whose existence is (F1.a)),
this gives the equivalence \(\Phi_{S}^{\mathrm{loc}}\).

\textbf{Step 4: Tensor compatibility with Hecke.} The tensor product
of \(\mathcal{H}_{v}\)-equivariant functors is
\(\mathcal{H}_{S}:=\otimes_{v\in S}\mathcal{H}_{v}\)-equivariant in the
appropriate sense. Hecke-functoriality across places follows from
cross-place commutativity
(\Cref{v4-arith:acd:prop:cross-place-commutativity}).

\textbf{Step 5: Spectral side decomposition.} The spectral side
\(\mathrm{IndCoh}_{\mathrm{Nilp}}(\prod_{v\in S}\mathcal{X}^{\mathrm{EG},v}_{GL_{2}})\)
decomposes as a tensor product of local factors because
\(\prod_{v\in S}\mathcal{X}^{\mathrm{EG},v}_{GL_{2}}\) is a product of
derived formal stacks (Gaitsgory--Rozenblyum's AMS DAG monographs give
continuity of \(\mathrm{IndCoh}_{\mathrm{Nilp}}\) under products of
derived stacks). Hence the right-hand side
of~\eqref{v4-arith:acd:eq:glued-FS} agrees with
\(\bigotimes_{v\in S}\mathrm{IndCoh}_{\mathrm{Nilp}}(\mathcal{X}^{\mathrm{EG},v}_{GL_{2}})\).

\textbf{Step 6: Whittaker-model compatibility.} The cross-place
compatibility requires checking that the Whittaker model on the
automorphic side, which factorizes as a tensor product over places,
matches the corresponding factorization on the spectral side. This
factorization is where (BC-diamond-glue) enters: the Whittaker-model
factorization is a compatibility between local Bernstein centres,
exactly what (BC-diamond-glue) asserts via compatibility~(b) of
\Cref{v4-arith:acd:thm:final-reduction}.

Thus \((L_{2})\) is not closed here. It is reduced to the local
categorical comparison (F1.a) plus (BC-diamond-glue).
\end{proof}

\begin{remark}[finite-prime \((L_{2})\) is conditional]
\label{v4-arith:acd:rem:L2-finite-unconditional}
For the subset of conditions involving only finite primes (no
archimedean place), the tensor-product formalism is unconditional
once the local categorical comparisons are supplied. Fargues--Scholze
arXiv:2102.13459, Lurie HA 4.8.1, and Gaitsgory--Rozenblyum supply
the local geometry, tensor functoriality, and product-continuity of
\(\mathrm{IndCoh}_{\mathrm{Nilp}}\). They do not prove the local
comparison with the Emerton--Gee stack, nor the Whittaker
compatibility encoded by (BC-diamond-glue).
\end{remark}

\section{Subcondition \((L_{3})\): Bernstein Centre Compatibility}
\label{v4-arith:acd:L3}

\subsection{Statement and Residual Content}

\begin{theorem}[\((L_{3})\) is (BC-diamond-glue)]
\label{v4-arith:acd:thm:L3-is-bc-diamond-glue}
\ClaimStatusProvedHere. \((L_{3})\) is equivalent to the
(BC-diamond-glue) compatibility of
\Cref{v4-arith:acd:thm:final-reduction}.
\end{theorem}

\begin{proof}
Forward: \((L_{3})\) asserts that the local Bernstein-centre actions
assemble to a global action on \(\bigotimes'_{v}\mathcal{C}_{v}\)
compatible with restricted tensor. The global Bernstein centre
\(\mathcal{Z}_{\mathrm{glob}}\) must embed into the colimit
\(\varinjlim_{S}\otimes_{v\in S}\mathcal{Z}_{v}\); this is precisely
compatibility~(c) of (BC-diamond-glue).

Reverse: given (BC-diamond-glue), conditions (a), (b), (c) supply
exactly the data required to construct the global action
\(\mathcal{Z}_{\mathrm{glob}}\curvearrowright\bigotimes'_{v}\mathcal{C}_{v}\)
as a filtered colimit of finite-\(S\) actions, with the compatibility
required by \((L_{3})\).
\end{proof}

\subsection{Partial change via Coherent Springer Theory}

\begin{theorem}[Ben-Zvi--Chen--Helm--Nadler arXiv:2010.02321]
\label{v4-arith:acd:thm:BZCHN}
\ClaimStatusProvedElsewhere. The local Bernstein centre
\(\mathcal{Z}_{p}\) of the categorical Hecke algebra at a prime \(p\)
embeds as a commutative dg-algebra into the global sections of a
coherent sheaf \(\mathcal{F}_{p}\) on the local Emerton--Gee stack
\(\mathcal{X}^{\mathrm{EG},p}_{GL_{2}}\) (``coherent Springer
sheaf''). Hecke operators at \(p\) correspond to action of
\(\mathcal{F}_{p}\) under the ind-coherent-sheaf correspondence.
\end{theorem}

\begin{remark}[what BZCHN supplies]
\label{v4-arith:acd:rem:BZCHN-supplies}
BZCHN supplies the local Bernstein centres as coherent sheaves on
local Emerton--Gee stacks. (BC-diamond-glue) asks that these
coherent sheaves glue across primes to a coherent sheaf on the
Selmer-derived global spectral stack. Local-to-global gluing is a standard
question in algebraic geometry: under standard flatness hypotheses,
coherent sheaves can descend along the Selmer-derived fibre-product
diagram. The non-standard content is \emph{agreement} of the
glued sheaf with the global Bernstein-centre action, a question
about the global automorphic category's Bernstein decomposition
(Bernstein 1984, categorified in Raskin 2021 JAMS and AGKRRV
2020--2024).
\end{remark}

\begin{conjecture}[BC-diamond-glue, concrete form]
\label{v4-arith:acd:cj:bc-diamond-glue-concrete}
The coherent Springer sheaves \(\{\mathcal{F}_{v}\}_{v}\) of
Ben-Zvi--Chen--Helm--Nadler (\Cref{v4-arith:acd:thm:BZCHN}) glue
along the diagonal of the Selmer-derived global spectral stack to a global
coherent sheaf \(\mathcal{F}_{\mathrm{glob}}\) whose action on
\(\mathrm{IndCoh}_{\mathrm{Nilp}}(\mathcal{X}^{\mathrm{Sel}}_{\bar\rho,\det,\mathcal{L}})\)
agrees, under the conjectural global arithmetic AG equivalence, with
the action of the global Bernstein centre on
\(\mathrm{D}([GL_{2}])^{\mathrm{sph\text{-}fin}}\).
\end{conjecture}

\begin{remark}[domain of the residual]
\label{v4-arith:acd:rem:L3-residual-domain}
\Cref{v4-arith:acd:cj:bc-diamond-glue-concrete} is the residual
content of \(\mathrm{L}_{\mathrm{ACD}}\) after the reduction of
\Cref{v4-arith:acd:thm:final-reduction}. Its domain:
\begin{itemize}
\item Gluing of coherent sheaves across a product of stacks:
non-formal theorem, standard techniques.
\item Agreement with global Bernstein centre under conjectural
global AG: non-formal theorem, requires careful adaptation of
the geometric-Langlands Bernstein-centre construction (AGKRRV) to the
number-field setting.
\item Total condition of (BC-diamond-glue): non-formal theorem.
\end{itemize}
This is \(1/3\) to \(1/2\) the original non-formal theorem of
\(\mathrm{L}_{\mathrm{ACD}}\).
\end{remark}

\section{Subcondition \((L_{4})\): Archimedean Extension}
\label{v4-arith:acd:L4}

\subsection{Statement}

\begin{theorem}[\((L_{4})\) archimedean boundary datum]
\label{v4-arith:acd:thm:L4-closed}
\ClaimStatusProvedHere. The archimedean factor
\(\mathcal{C}_{\infty}=\mathrm{D}^{\mathrm{sph\text{-}fin}}([GL_{2}(\mathbb{R})])\)
is modelled by the ind-completed derived category of admissible
\((\mathfrak{gl}_{2},O(2))\)-modules of finite \(K\)-type.  It carries
the trivial \(K\)-spherical object as distinguished unit and extends
the finite-place restricted tensor product by direct tensoring.
\end{theorem}

\begin{proof}
\textbf{Step 1: Well-definedness.} The derived category of admissible
\((\mathfrak{g}, K_{\infty})\)-modules, where
\(\mathfrak{g}=\mathfrak{gl}_{2}(\mathbb{R})\) and
\(K_{\infty}=O(2)\subseteq GL_{2}(\mathbb{R})\), has the standard
Harish-Chandra description (Jacquet--Langlands SLN 114~\S1--3;
Vogan 1981 ``Representations of Real Reductive Groups''). Its
ind-completion is stable and presentable. The sph-fin cut restricts to
finite \(K_{\infty}\)-type objects.

\textbf{Step 2: Distinguished unit.} The spherical unit at
\(v=\infty\) is the trivial \(K_{\infty}\)-spherical object.  Spherical
principal series are not unique; they are objects of the category, not
the distinguished unit.

\textbf{Step 3: Tensor product extension.} The restricted tensor
product of \Cref{v4-arith:acd:def:restricted-tensor} extends
naturally by including \(v=\infty\) in the index set \(V\) of places:
\(V=\mathbb{P}\cup\{\infty\}\) with cofinite unramified subset
\(V^{\mathrm{unr}}=\mathbb{P}\cup\{\infty\}\setminus S_{\mathrm{ram}}\)
for the given ramification set \(S_{\mathrm{ram}}\). The archimedean
factor enters as a tensor factor in each \(S\) containing \(\infty\) and
by identity-via-distinguished-unit in \(S\) excluding \(\infty\).

\textbf{Step 4: Hecke action at infinity.} The archimedean Hecke
action is represented by the centre of the universal enveloping algebra
\(U(\mathfrak{gl}_{2}(\mathbb{R}))\), commuting with finite-place
Hecke action by construction (both act via representation-theoretic
restriction to the adelic double coset).

\textbf{Step 5: Archimedean local factors.} Jacquet--Langlands SLN
114 \S1--3 classify the archimedean local representations of
\(GL_{2}(\mathbb{R})\) and \(GL_{2}(\mathbb{C})\). The sph-fin
subcategory is finite-type in each isotypic component; Vogan's
classification gives a complete parametrization by
Harish-Chandra parameters.

This proves the archimedean boundary datum. It does not construct a
number-field spectral equivalence at \(v=\infty\).
\end{proof}

\begin{remark}[archimedean spectral side]
\label{v4-arith:acd:rem:archimedean-spectral}
The spectral-side analogue at \(v=\infty\) is
a proposed Weil-parameter stack for
\(W_{\mathbb{R}}=\mathbb{C}^{\times}\rtimes\mathbb{Z}/2\).  Vogan and
Langlands classify the representations.  They do not give an
Emerton--Gee stack at infinity, and no local categorical equivalence at
infinity is asserted here.
\end{remark}

\section{Curve Analogy: Obstructions and Transfers}
\label{v4-arith:acd:function-field-transport}

\subsection{The Curve Theorem}

\begin{theorem}[de Rham geometric Langlands, AGKRRV]
\label{v4-arith:acd:thm:FF-geom-langlands}
\ClaimStatusProvedElsewhere. For a smooth complete curve \(X\) over an
algebraically closed field of characteristic zero and a reductive group
\(G\),
there is an equivalence
\begin{equation}
\mathrm{D}\text{-}\mathrm{mod}(\mathrm{Bun}_{G}(X))\;\simeq\;\mathrm{IndCoh}_{\mathrm{Nilp}}(\mathrm{LocSys}^{\mathrm{dR}}_{G}(X))
\end{equation}
satisfying Hecke-functoriality, Bernstein-centre compatibility, and
categorical descent over the curve. The proof combines
Gaitsgory--Rozenblyum's AMS DAG monographs, Raskin's chiral principal
series and D-module foundations, and the AGKRRV mathematical treatments on geometric
Langlands.
\end{theorem}

\subsection{Obstructions to Transport}

\begin{theorem}[obstructions to number-field transport]
\label{v4-arith:acd:thm:transport-obstructions}
\ClaimStatusProvedHere. The de Rham geometric-Langlands proof of
\Cref{v4-arith:acd:thm:FF-geom-langlands} does not transport directly
to the number-field case. The obstructions:
\begin{description}
\item[(Obs.1)] \emph{No archimedean place.} Geometric Langlands on a
curve has no real place. \((L_{4})\) has no curve analogue
and requires separate treatment (Jacquet--Langlands archimedean).
\item[(Obs.2)] \emph{Non-smooth Emerton--Gee stack.} Number-field
\(\mathcal{X}^{\mathrm{EG}}_{GL_{2}}\) is formally algebraic but not
smooth over \(\mathrm{Spec}\,\mathbb{Z}\); the de Rham stack
\(\mathrm{LocSys}^{\mathrm{dR}}_{G}(X)\) is smooth in the standard
geometric-Langlands range
under mild hypotheses. This affects the Bernstein-centre
compatibility \((L_{3})\).
\item[(Obs.3)] \emph{Perfectoid diamonds vs ind-schemes.}
Geometric Langlands uses algebraic and ind-algebraic loop spaces.
Number-field local geometrisation uses perfectoid diamonds (Scholze
arXiv:1709.07343). The gluing formalism differs, and the
Bhatt--Lurie prismatic comparison arXiv:2201.06120 is required to
bridge the two.
\item[(Obs.4)] \emph{Integrality of Galois deformations.}
De Rham local systems are geometric; number-field Galois
deformations are integral (Emerton--Gee stack). The Bernstein-centre
action incorporates integrality in a way the de Rham curve case
does not.
\end{description}
\end{theorem}

\begin{corollary}[what the curve analogy supplies]
\label{v4-arith:acd:cor:transport-summary}
\ClaimStatusProvedHere. Under the de Rham curve proof, the
following patterns survive in the number-field setting:
\begin{itemize}
\item \((L_{1})\) (restricted tensor product well-definedness):
\textbf{formal pattern}. Raskin's construction is
\(\infty\)-categorical and does not use curve-specific structure beyond
existence of local categories with distinguished units.
\item \((L_{2})\) (Hecke-functoriality gluing over a finite set of
primes): \textbf{conditional pattern}. It applies only after the local
categorical comparison is supplied. Fargues--Scholze replaces the local de~Rham
geometric input by diamond geometry and spectral action; the
cross-place compatibility argument still needs the local comparison
and (BC-diamond-glue).
\end{itemize}
The following do \emph{not} transport:
\begin{itemize}
\item \((L_{3})\) (Bernstein-centre compatibility): obstructed by
(Obs.2), (Obs.4). Resolution requires (BC-diamond-glue),
\Cref{v4-arith:acd:cj:bc-diamond-glue-concrete}.
\item \((L_{4})\) (archimedean extension): obstructed by (Obs.1).
Resolution requires independent Jacquet--Langlands archimedean
construction (\S\ref{v4-arith:acd:L4}).
\end{itemize}
\end{corollary}

\section{Final Reduction: \(\mathrm{L}_{\mathrm{ACD}}\) to a Single Lemma}
\label{v4-arith:acd:final-reduction}

\subsection{Statement}

\begin{theorem}[\(\mathrm{L}_{\mathrm{ACD}}\) residual reduction]
\label{v4-arith:acd:thm:LACD-single-lemma}
\ClaimStatusNeedsVerification. \(\mathrm{L}_{\mathrm{ACD}}\) holds after
the local categorical comparison (F1.a) and the BC-diamond-gluing
compatibility lemma
(\Cref{v4-arith:acd:cj:bc-diamond-glue-concrete}) are supplied.
\end{theorem}

\begin{proof}
By \Cref{v4-arith:acd:thm:modular-decomposition},
\(\mathrm{L}_{\mathrm{ACD}}\) contains
\((L_{1})\), \((L_{2})\), \((L_{3})\), and \((L_{4})\), together with
the global identification.

\((L_{1})\) is a theorem (\Cref{v4-arith:acd:cor:L1-holds}).

\((L_{2})\) is conditional: it needs the local categorical
comparison (F1.a) and (BC-diamond-glue) for the cross-place
Whittaker compatibility
(\Cref{v4-arith:acd:thm:L2-reduction}, Steps 1 and 6).

\((L_{3})\) is controlled by the Bernstein-centre gluing datum
(\Cref{v4-arith:acd:thm:L3-is-bc-diamond-glue}).

\((L_{4})\) is the archimedean boundary datum
(\Cref{v4-arith:acd:thm:L4-closed}).

Hence \(\mathrm{L}_{\mathrm{ACD}}\) reduces, in this chapter, to the
conjunction of F1.a, (BC-diamond-glue), and the global identification
after the formal construction of \((L_{1})\) and the archimedean datum
\((L_{4})\).
\end{proof}

\subsection{Dependency on Established Theorems}

\begin{proposition}[dependency table of \(\mathrm{L}_{\mathrm{ACD}}\)]
\label{v4-arith:acd:prop:dependency-closure}
\ClaimStatusProvedHere. The conditional reduction of
\(\mathrm{L}_{\mathrm{ACD}}\) uses the following established theorems:
\begin{enumerate}
\item Fargues--Scholze arXiv:2102.13459 (local geometrization at each
prime).
\item Emerton--Gee arXiv:1908.07185 (local Emerton--Gee stack
construction; the global target is Selmer-derived here).
\item Gaitsgory--Rozenblyum, \emph{A Study in Derived Algebraic
Geometry}, AMS Mathematical Surveys and Monographs 221, Vols.~I--II
(\((\infty,2)\)-category of correspondences,
\(\mathrm{IndCoh}_{\mathrm{Nilp}}\) continuity).
\item Raskin, \emph{Chiral principal series categories I}, Adv.
Math. 388 (2021), 107856, and \emph{D-modules on infinite-dimensional
varieties} (restricted tensor product machinery for principal-series
categories).
\item Lurie HA 4.8 (tensor products of presentable \(\infty\)-categories),
5.3 (compact generation, higher Deligne).
\item Scholze arXiv:1709.07343 (\'etale cohomology of diamonds).
\item Bhatt--Scholze arXiv:1905.08229 (prismatic cohomology).
\item Bhatt--Lurie arXiv:2201.06120 (absolute prismatic,
Nygaard-filtration convention bridge).
\item Zhu arXiv:1603.05593 (Witt-vector affine Grassmannian).
\item Ben-Zvi--Chen--Helm--Nadler arXiv:2010.02321 (coherent Springer
theory; local Bernstein centres as coherent sheaves).
\item Jacquet--Langlands SLN 114 (archimedean local factors).
\item Vogan 1981 ``Representations of Real Reductive Groups''
(archimedean classification).
\end{enumerate}
Together with the local categorical comparison (F1.a),
(BC-diamond-glue), and the global identification, these established
results constitute the primary literature input for the reduction.
\end{proposition}

\begin{remark}[independent decomposition]
\label{v4-arith:acd:rem:source-disjoint-decomposition}
The twelve established theorems of
\Cref{v4-arith:acd:prop:dependency-closure} decompose into the
following independent source lists:
\begin{description}
\item[Fargues--Scholze list:] (1), (6), (7). Perfectoid / diamond
geometry, Fargues--Fontaine curve.
\item[Gaitsgory--Rozenblyum list:] (3), (4). DAG and
\((\infty,2)\)-categorical correspondences, restricted tensor.
\item[Lurie list:] (5). Higher algebra foundations.
\item[Emerton--Gee list:] (2). Galois deformation stacks.
\item[Bhatt--Lurie list:] (8). Prismatic Nygaard bridge.
\item[Zhu list:] (9). Witt-vector Grassmannian.
\item[Coherent Springer theory list:] (10). Local Bernstein
centres as coherent sheaves.
\item[Jacquet--Langlands + Vogan list:] (11), (12). Archimedean
local representation theory.
\end{description}
These eight lists are disjoint in the orthogonal sense: no pair shares
a primary author at overlapping work, and each list has an
independent primary literature. A future proof of
\(\mathrm{L}_{\mathrm{ACD}}\) through this construction must keep the same
source separation.
\end{remark}

\section{Consequence: F1 Closure Map}
\label{v4-arith:acd:F1-closure-map}

\subsection{The Full Chain}

\begin{theorem}[F1 residual chain via \(\mathrm{L}_{\mathrm{ACD}}\)]
\label{v4-arith:acd:thm:F1-closure}
\ClaimStatusNeedsVerification. F1 (global arithmetic
Arinkin--Gaitsgory, \Cref{v4-arith:conj:F1-global-AG}) reduces to
the local categorical comparison (F1.a) plus
(BC-diamond-glue) (\Cref{v4-arith:acd:cj:bc-diamond-glue-concrete}).
\end{theorem}

\begin{proof}
By \Cref{v4-arith:cl:cor:F1-single-lemma}, F1 reduces to F1.a plus
\(\mathrm{L}_{\mathrm{ACD}}\). By
\Cref{v4-arith:acd:thm:LACD-single-lemma},
\(\mathrm{L}_{\mathrm{ACD}}\) reduces to F1.a plus (BC-diamond-glue).
The theorem identifies the residual pair \{F1.a, (BC-diamond-glue)\}
as the obstruction visible in this reduction.
\end{proof}

\subsection{Formal Consequences}

\begin{corollary}[consequent form]
\label{v4-arith:acd:cor:consequent-form}
\ClaimStatusNeedsVerification. Given the local categorical comparison
(F1.a), (BC-diamond-glue), F2, and the corrected arithmetic BZN
comparison, the following formal consequences become eligible for
theorem-grade closure:
\begin{itemize}
\item F1 (global arithmetic Arinkin--Gaitsgory).
\item F3.BZN\(^{\mathrm{Ar}}\) (arithmetic Ben-Zvi--Nadler, conditional
on F1 \(+\) F2 plus the full-target/formal-loop comparison).
\item The six-window placement theorem
(\Cref{v4-arith:thm:six-window-placement}) remains conditional on
the same residual chain.
\end{itemize}
\end{corollary}

\begin{remark}[six-window upgrade]
\label{v4-arith:acd:rem:six-window-upgrade}
The upgrade of \Cref{v4-arith:thm:six-window-placement} from
\ClaimStatusConjectured to \ClaimStatusProvedHere requires F1 (hence
(BC-diamond-glue)), F2 (hence the \(\infty\)-categorical Tate
restricted tensor lemma F2.c), and F3 (hence F3.BZN\(^{\mathrm{Ar}}\)).
With (BC-diamond-glue) and the local comparison supplied, F1 becomes a
conditional theorem.  With F2.c supplied, F2 becomes a conditional
theorem.  F3 then requires F3.BZN\(^{\mathrm{Ar}}\). The six-window
placement is theorem-grade only on the conjunction of these residual
inputs.
\end{remark}

\subsection{Effect on \(\mathcal{V}^{\mathrm{prim}}\) Spectral Identification}

\begin{corollary}[\(\mathcal{V}^{\mathrm{prim}}\) spectral shadow]
\label{v4-arith:acd:cor:V-prim-spectral}
\ClaimStatusNeedsVerification. Given the local categorical comparison,
(BC-diamond-glue), and the corrected arithmetic BZN full-target
comparison, the spectral shadow of \(\mathcal{V}^{\mathrm{prim}}\) is
\begin{equation}
\mathcal{V}^{\mathrm{prim}}\;\rightsquigarrow\;
\mathcal{O}\bigl(L\mathcal{X}^{\mathrm{Sel}}_{\bar\rho,\det,\mathcal{L}}\bigr).
\end{equation}
The further passage to
\(\mathcal{O}(\mathcal{X}^{\mathrm{Sel}}_{\bar\rho,\det,\mathcal{L}})\)
requires a constant-loop restriction or a flat-locus truncation. It is
not a consequence of Hochschild trace alone.
\end{corollary}

\begin{proof}
\(\mathcal{V}^{\mathrm{prim}}=\mathrm{HH}_{\bullet}(\mathrm{Hk}^{\mathrm{Iw}}_{GL_{2},\mathrm{glob}})\)
(\Cref{v4-arith:thm:crowning-placement}). Under F1, which here means
the local comparison together with (BC-diamond-glue),
\(\mathrm{Hk}^{\mathrm{Iw}}_{GL_{2},\mathrm{glob}}\simeq\mathrm{IndCoh}_{\mathrm{Nilp}}(\mathcal{X}^{\mathrm{Sel}}_{\bar\rho,\det,\mathcal{L}})\).
Ben-Zvi--Nadler arXiv:0904.1247 computes
\(\mathrm{HH}_{\bullet}(\mathrm{IndCoh}_{\mathrm{Nilp}}(\mathcal{X}))=\mathcal{O}(L\mathcal{X})=\mathcal{O}(\mathrm{Map}(S^{1},\mathcal{X}))\)
for derived stacks \(\mathcal{X}\). The constant-loop map
\(\mathcal{X}\to L\mathcal{X}\) gives a restriction from the loop
shadow to \(\mathcal{O}(\mathcal{X})\).  This restriction is additional
data.  The categorical trace supplies the loop object.
\end{proof}

\section{Residual Obstructions}
\label{v4-arith:acd:residuals}

\subsection{Obstruction Decomposition}

After the reduction of \Cref{v4-arith:acd:thm:final-reduction}, the
residual inputs for \(\mathrm{L}_{\mathrm{ACD}}\), and hence for F1,
are:

\begin{itemize}
\item \textbf{Local categorical comparison:} the passage from the
pointwise local spectral action to the local categorical equivalence
used in \(L_{2}\).
\item \textbf{Global restricted-tensor identification:} the comparison
between the restricted tensor product and the global automorphic
category.
\item \textbf{(BC-diamond-glue):}
\Cref{v4-arith:acd:cj:bc-diamond-glue-concrete}. The central open
residual, non-formal theorem. Domain: gluing of coherent Springer
sheaves across a product of local Emerton--Gee stacks, plus
agreement with the global Bernstein centre acting on the global
automorphic category.
\end{itemize}

\subsection{Prior Hypotheses}

(BC-diamond-glue) is independent of the remaining open items from
the arithmetic branch of Vol~IV:
\begin{itemize}
\item F2.c (F2 residual, \(\infty\)-categorical Tate tensor): independent branch.
\item Arithmetic Positselski: independent, on the P1 branch.
\item Koszul--Petersson compatibility: independent, on the P3 branch.
\item VB\(^{*}\) (arithmetic Virasoro positivity): independent, on the
RH branch.
\end{itemize}

(BC-diamond-glue) is central on the F1 branch only after the local
comparison and the global identification are fixed.

\subsection{Primary-Literature Precedents}

\begin{itemize}
\item Ben-Zvi--Chen--Helm--Nadler arXiv:2010.02321: \emph{local}
Bernstein centres as coherent sheaves on local Emerton--Gee stacks.
This is the state of the art; (BC-diamond-glue) is the global
assembly of these local constructions.
\item Gan--Savin: partial change on restricted
tensor product of local automorphic categories, with specific
compatibilities verified for certain Jacquet--Langlands transfers;
the global Bernstein-centre compatibility for \(GL_{2}/\mathbb{Q}\)
remains open.
\item AGKRRV 2020--2024 de Rham geometric Langlands: supplies the
curve analogue of the centre-compatibility pattern.  It does not prove
the number-field (BC-diamond-glue) statement. The surviving pattern is
documented in
\Cref{v4-arith:acd:cor:transport-summary}.
\end{itemize}

\subsection{Paper-Length Estimate, Revised}

\begin{table}[h]
\centering
\small
\begin{tabular}{l|l|l|l}
Residual & Original condition & Revised condition & Form \\
\hline
\(\mathrm{L}_{\mathrm{ACD}}\) & non-formal theorem & conditional & open \\
\((L_{1})\) & n/a & closed (closed) & Theorem, \Cref{v4-arith:acd:cor:L1-holds} \\
\((L_{2})\) & n/a & non-formal theorem plus local comparison & conditional \\
\((L_{3})\) = (BC-diamond-glue) & n/a & non-formal theorem & Conjectural residual \\
\((L_{4})\) & n/a & boundary datum & no spectral equivalence \\
\end{tabular}
\caption{Revised substantial estimate for \(\mathrm{L}_{\mathrm{ACD}}\)
closure after the reduction of \Cref{v4-arith:acd:thm:final-reduction}.}
\label{v4-arith:acd:table:revised-domain}
\end{table}

Total revised domain for \(\mathrm{L}_{\mathrm{ACD}}\) closure:
non-formal theorem plus the local categorical comparison and the
global identification.  The number records the visible
Bernstein-centre and Whittaker work; it is not a proof that all
adelic descent has been closed.

\section{Descent Boundary}
\label{v4-arith:acd:summary}

\begin{enumerate}
\item \(\mathrm{L}_{\mathrm{ACD}}\) decomposes into four disjoint
subconditions \((L_{1})\)--\((L_{4})\)
(\Cref{v4-arith:acd:thm:modular-decomposition}).
\item \((L_{1})\) (restricted tensor well-definedness) closes
unconditionally via Lurie HA 4.8 \(+\) 5.3, Raskin 2021--2022
(\Cref{v4-arith:acd:thm:L1-closed},
\Cref{v4-arith:acd:cor:L1-holds}).
\item \((L_{4})\) (archimedean extension) supplies the
representation-theoretic boundary datum via Jacquet--Langlands SLN 114
\(+\) Vogan 1981
(\Cref{v4-arith:acd:thm:L4-closed}).
\item \((L_{2})\) (Hecke-functoriality gluing) is conditional:
after the local categorical comparisons are supplied it reduces
to (BC-diamond-glue), via tensor product formalism \(+\)
Gaitsgory--Rozenblyum continuity
(\Cref{v4-arith:acd:thm:L2-reduction}).
\item \((L_{3})\) is the Bernstein-centre part of (BC-diamond-glue)
(\Cref{v4-arith:acd:thm:L3-is-bc-diamond-glue}).
\item \(\mathrm{L}_{\mathrm{ACD}}\) reduces to local comparison,
(BC-diamond-glue), the archimedean boundary datum, and the global
identification
(\Cref{v4-arith:acd:thm:LACD-single-lemma}).
\item F1 (global arithmetic Arinkin--Gaitsgory) is conditional on the
same residual chain
(\Cref{v4-arith:acd:thm:F1-closure}).
\item Given the residual chain, the six-window placement theorem
is conditional theorem-grade.  The Hochschild trace of
\(\mathcal{V}^{\mathrm{prim}}\) sees the loop stack; passage to the
Selmer-derived structure sheaf is an extra constant-loop or flat-locus
restriction
(\Cref{v4-arith:acd:cor:V-prim-spectral}).
\item Revised condition for closure: non-formal theorem, reduced from
non-formal theorem (\Cref{v4-arith:acd:table:revised-domain}).
\end{enumerate}

\(\mathrm{L}_{\mathrm{ACD}}\) reduces from one large conjecture to a
formal tensor construction, local comparison, Bernstein-centre gluing,
an archimedean boundary datum, and a global identification.  The
structural content is a reduction, not an equivalence theorem for F1.
The reduction is supported by eight disjoint primary-literature sources
(\Cref{v4-arith:acd:rem:source-disjoint-decomposition}).

The residual (BC-diamond-glue) is a single compatibility check
between Ben-Zvi--Chen--Helm--Nadler's local coherent Springer
sheaves and the global Bernstein centre on the global automorphic
category. It is central on the F1 six-window branch of the
arithmetic branch, but it is not the only hypothesis.

\section{Forward Pointers}
\label{v4-arith:acd:forward}

\begin{itemize}
\item Chapter~\ref{v4-arith:closure} gives the parent reduction of F1
(\Cref{v4-arith:thm:F1-reduction}) into sub-items (F1.a)--(F1.d).
\Cref{v4-arith:acd:thm:modular-decomposition} refines that reduction
at the adelic-descent layer.
\item Chapter~\ref{v4-arith:synthesis},
\Cref{v4-arith:thm:six-window-placement}, is upgraded (conditional
on (BC-diamond-glue) and the F2 / F3 residuals) to theorem-grade via
\Cref{v4-arith:acd:cor:consequent-form}.
\item Chapter~\ref{v4-arith:rh-reformulation} contains the RH
reformulation F-RH; (BC-diamond-glue) sits on the six-window-placement
branch (F1--F3 structural branch in
Appendix~\ref{v4-arith:app:three-conjectures-dep}), not the
RH-closure branch (four-input RH branch). (BC-diamond-glue) controls the
categorical-trace presentation of \(\mathcal{V}^{\mathrm{prim}}\);
the RH-closure branch (four-input RH branch)
remains on a separate dependency chain.
\item Appendix~\ref{v4-arith:app:three-conjectures-dep},
\Cref{v4-arith:app:C:residual-list}, records
\(\mathrm{L}_{\mathrm{ACD}}\) on the F1 branch.  After the reduction of
\Cref{v4-arith:acd:thm:final-reduction}, the Bernstein-centre part is
(BC-diamond-glue), with visible requires a non-formal theorem.
\end{itemize}
